% Paper 6 (GRPO-Registry) schema section
\section{Registry Schema}
\label{sec:p6-schema}

The schema (\texttt{registry/schema.json}) defines two record types under one
\texttt{oneOf}. Both carry a \texttt{schema\_version} for forward migration
(Section~\ref{sec:p6-governance}). A single field convention governs
everything: \textbf{JSON \texttt{null} means \emph{unreported}}, while explicit
\texttt{false} / \texttt{0.0} / \texttt{"none"} means
\emph{reported-as-absent}. A stack that says ``we use no KL term'' is making a
checkable claim; a stack that says nothing is not. The auditor badge
(Section~\ref{sec:p6-query}) scores exactly this distinction: reporting
coverage, not configuration virtue.

\subsection{Stack Records: the Seven-Field MIN-REPORT Run Manifest}
\label{sec:p6-schema-stack}
A \texttt{stack} record describes one deployed GRPO-family training stack. Its
core is the \texttt{min_report} object, one sub-object per item of the
MIN-REPORT-RL minimum reportable stack:

\begin{enumerate}[leftmargin=1.8em, itemsep=0.1em]
  \item \textbf{loss_form}: importance-ratio level (token/sequence), clip
        bounds ($\eps_{\text{low}}, \eps_{\text{high}}$), length
        normalization, advantage normalization (std/none), token mask. These
        are the fields the variant papers fight over.
  \item \textbf{reference_kl}: whether a reference policy exists, the KL
        coefficient $\beta$, and the estimator/surrogate used.
  \item \textbf{sampler_backend}: backend identity, numeric precision, and sampler
        parameters (temperature, top-$p$). The \textbf{base-checkpoint
        identity (revision/hash)} this item requires lives in the top-level
        \texttt{model} field and \texttt{run\_config} \texttt{checkpoint.*};
        \texttt{stackdiff} compares it at R5, but it is not a badge leaf. Our
        nominal backend comparison ($85.6\%$ vs.\ $5.0\%$) used Base vs.\
        Instruct checkpoints under one label, so a record without it
        cannot distinguish a framework effect from a model swap.
  \item \textbf{telemetry}: whether the stack emits a \emph{per-step} ZVF and
        gradient-utilization ($\mathrm{GU}=1-\mathrm{ZVF}$) trajectory, and
        from what source. Endpoint summaries are not enough: ZVF is high both
        when a policy has mastered a prompt and when it cannot solve it
        ($h_G(p)=h_G(1-p)$), so a single mean aliases mastery with incapacity
        (see the companion ZVF paper).
  \item \textbf{group_size_schedule}: initial $G$, fixed vs.\ adaptive, and
        the adaptation rule if any.
  \item \textbf{heldout_split}: whether evaluation prompts are disjoint from
        the reward environment, and how the split was built.
  \item \textbf{decontamination}: decontamination performed, and whether a
        parser-robustness probe was run against the reward extractor.
\end{enumerate}

This seven-field object is the run-manifest portion of the eight-item standard;
item~8 (held-out pass@$k$ curves) belongs to the evaluation report and may be
attached under outcomes rather than being known at run start.

Around the \texttt{min_report} block sit provenance
(\texttt{source_artifacts}, W\&B project, date), identity
(\texttt{label_claimed}, framework name/version, and an
\texttt{openness} enum: \texttt{open}/\texttt{managed}/\texttt{closed}), an
optional \texttt{outcomes} summary (informational only---never part of the
badge), and \texttt{variant_deltas_applied}, described next.

\subsection{Variant-Delta Records}
\label{sec:p6-schema-delta}
A \texttt{variant_delta} record decomposes a named variant into components,
each mapping a \texttt{min_report} leaf (e.g.\
\texttt{loss_form.clip_eps_high}) to a change relative to base GRPO. We ship
18. Three are verified component by component against their source papers
rather than against folklore; the other 15 are not (four have no
peer-reviewed source, and AERO's citation describes a different mechanism):

\begin{itemize}[leftmargin=1.5em, itemsep=0.15em]
  \item \textbf{DAPO}~\citep{yu2025dapo}: five components---%
        \texttt{clip_higher} (asymmetric/decoupled clip,
        $\eps_{\text{low}}{=}0.2$, $\eps_{\text{high}}{=}0.28$),
        \texttt{dynamic_sampling} (filter zero-variance groups and resample
        until the batch is full), \texttt{token_level_loss} (token-level
        policy-gradient aggregation), \texttt{overlong_reward_shaping} (soft
        length-aware penalty), and \texttt{kl_removed} (the KL term is
        excluded from the objective).
  \item \textbf{Dr.~GRPO}~\citep{tong2025drgrpo}: two removals---the
        $1/|o_i|$ response-length normalization in the loss, and the division
        by the within-group reward standard deviation in the advantage.
  \item \textbf{GSPO}~\citep{qwen2025gspo}: the importance ratio is defined on
        (length-normalized) sequence likelihood rather than per token, with
        clipping and optimization applied at the sequence level.
\end{itemize}

\subsection{Where Label Flips Become Machine-Readable}
\label{sec:p6-schema-status}
The joint between the two record types is the stack-side
\texttt{variant_deltas_applied} array: for each component of each claimed
delta, the stack declares a status in
$\{\texttt{implemented}, \texttt{surrogate}, \texttt{absent},
\texttt{unknown}\}$. The shipped ``DAPO-on-managed-runtime'' entry, for
example, records \texttt{clip_higher: absent} (the arm clips the advantage;
its loss has no importance ratio to clip), \texttt{dynamic_sampling: absent},
\texttt{token_level_loss: absent}, \texttt{overlong_reward_shaping: absent},
and \texttt{kl_removed: implemented}, each with a pointer into the client-side
run script:

\begin{lstlisting}[basicstyle=\ttfamily\scriptsize, frame=single,
  caption={Excerpt from \texttt{entries/tinker_dapo_qwen3.5-4b_gsm8k.json}
  (two of five components elided).},
  label={lst:p6-dapo-entry}]
"label_claimed": "dapo",
"variant_deltas_applied": [
  {"delta_id": "delta_dapo", "component": "clip_higher",
   "status": "absent",
   "note": "advantage clip [-2,3] (experiments/tinker-runs/
            cell_runner.py@b0ac85bfb:L181); no ratio exists to clip"},
  {"delta_id": "delta_dapo", "component": "dynamic_sampling",
   "status": "absent",
   "note": "no zero-variance filter/resample in the loop"},
  ...
  {"delta_id": "delta_dapo", "component": "kl_removed",
   "status": "implemented",
   "note": "custom loss has no KL term (experiments/tinker-runs/
            live_zvf_probe.py@b0ac85bfb:L78)"} ]
\end{lstlisting}

A reader---or a script---can now see at a glance that this stack's ``DAPO'' is
one component out of five, without reading a paper appendix or reverse
engineering a launcher.

\paragraph{Self-reported CI provenance.} The \texttt{outcomes} block carries
an optional \texttt{ci\_method} object (\texttt{method}, \texttt{n\_boot},
\texttt{seed}, \texttt{ci\_level}, \texttt{source}) so an entry can declare
\emph{how} the confidence interval on its reported telemetry was computed,
rather than leaving a reader to guess. The field is a backward-compatible
addition: it is optional and unconstrained by any \texttt{required} clause, so
existing entries continue to validate without modification. It is populated on
the $7$ managed-runtime Qwen3.5-4B entries: the three N2 same-stack entries
(Family~D) and the four head-to-head entries (Family~C). All seven name the
paired bootstrap of \texttt{registry\_validate.py} over aligned per-step N2
tensors as their source ($n_{\text{boot}}{=}2000$, seed $0$, percentile
$95\%$). The four Family~C entries come from the separate three-seed
head-to-head (seeds 42, 123, 456), and the DAPO, Dr.~GRPO, and GSPO arms have
no N2 tensors, so that source attribution cannot describe them; it is flagged
for correction in the shipped entries. The remaining $21$ of the
\NSTACK{} stack entries leave it \texttt{null} (unreported). This is an
informational disclosure gap: the outcome-blind MIN-REPORT badge excludes
\texttt{outcomes.ci\_method} and does not score its presence or absence.
Reproduce with \texttt{python3 scripts/p5p8/add\_ci\_method.py} (idempotent;
outputs \texttt{registry\_ci\_method\_coverage.tsv}).
